\section{Prerrequisito}

Practicas realizadas en entorno linux (Debian 12 - Bookworm)


\begin{itemize}
    
   \item \textbf{Eclipse Temurin JDK JAVA 25:} Entorno de desarrollo de Java (Java Development Kit) proporcionado por la fundación Eclipse Adoptium, utilizado para la compilación y ejecución de la aplicación backend.
    
   \item \textbf{DBeaver:} Cliente gráfico y gestor de bases de datos mulitplataforma, utilizado para la administración, ejecución de scripts DDL y DML, y visualización de la base de datos PostgreSQL.
   
   \item \textbf{IDE IntelliJ IDEA:} Entorno de desarrollo integrado utilizado para la escritura, depuración y gestión de la lógica de negocio y arquitectura en Java.
   
   \item \textbf{Git (Terminal / CLI):} Sistema de control de versiones distribuido que permite rastrear cambios en el código fuente, gestionar ramas de desarrollo y coordinar el historial del proyecto desde la línea de comandos.

   \item \textbf{Plataforma GitHub:} Servicio de alojamiento en la nube para repositorios Git, utilizado para el respaldo del código fuente, el control de versiones remoto y la integración del proyecto.
    
   \item \textbf{VSCodium:} Editor de código fuente basado en los archivos de Visual Studio Code, configurado desde su compilación como un binario de código abierto (\textit{Open Source}) y libre de telemetría, utilizado para la edición ligera de scripts, archivos de configuración y documentación.
   
   \item \textbf{Docker / Docker Compose:} Requisito indispensable para la orquestación y el levantamiento aislado de los servicios del backend y la base de datos PostgreSQL, garantizando la portabilidad del entorno.
    Para iniciar el contenedor de PostgreSQL en segundo plano y ver inmediatamente los registros (\textit{logs}) del servicio en tiempo real, se ejecuta la combinación de ambos comandos:
    \begin{verbatim}
        docker compose up -d
    \end{verbatim}
\end{itemize}


%% end \section{Prerrequisito}




\section{Requerimiento}


\subsection{Lista de requerimientos / Épicas}

    
   \begin{table}[H]
    \centering % Centra la tabla en la página
        \begin{tabular}{|c|c|c|} % Tres columnas centradas con líneas verticales
            \hline
            \textbf{ID}    &   \textbf{Acción}     &   \textbf{Ejemplo}        \\ \hline
            REQ-01         &   Avance tercera seccion  &   Desarrollo de la base de datos de la tercera seccion  \\ \hline
            
        \end{tabular}
    \caption{Comandos básicos de Git} % Título de la tabla
    \end{table}



    %%end  \subsection{Lista de requerimientos / Épicas}
    
%% end \section{Requerimiento}